\section{Simple tensors}

\begin{definition}[Physical-trace transfer of a doubled-index tensor]
    \label{def:mpdo_doubled_phys_trace_transfer}
    \lean{MPOTensor.doubledPhysTraceTransfer}
    \leanok
    \uses{def:mpo_tensor}
    For a doubled-index tensor $B$ with physical dimension $d^{2}$, the
    physical-trace transfer is the bond matrix obtained by closing the ket
    leg against the bra leg of one tensor:
    \begin{align}
        \mathcal T_{B}
         & =\sum_{i} B^{\,ii}.
        \label{eq:rfp_phys_trace}
    \end{align}
    This is the transfer object of the zero-correlation-length condition of
    \cite[Definition~4.2, lines~735--739]{Cirac2016MPDO_arXiv}, applied to
    one element of a basis of normal tensors.
\end{definition}

\begin{lemma}[Identification with the tensor-level transfer]
    \label{lem:mpdo_doubled_phys_trace_transfer_to_mps}
    \lean{MPOTensor.doubledPhysTraceTransfer_toMPSTensor}
    \leanok
    \uses{def:mpdo_doubled_phys_trace_transfer}
    The physical-trace transfer of the doubled-index view of an MPO tensor
    is the physical-trace transfer of the tensor itself.
\end{lemma}

\begin{proof}\leanok
    If $M^{\mathrm{dbl}}$ is the doubled-index view of $M$, then
    $(M^{\mathrm{dbl}})^{(i,j)}=M^{ij}$, and hence
    $\mathcal T_{M^{\mathrm{dbl}}}
        =\sum_i(M^{\mathrm{dbl}})^{ii}
        =\sum_i M^{ii}
        =\mathcal T_M$.
\end{proof}

\begin{lemma}[Blockwise form of the physical-trace transfer]
    \label{lem:mpdo_doubled_phys_trace_transfer_blockwise}
    \lean{MPOTensor.doubledPhysTraceTransfer_toTensor}
    \leanok
    \uses{def:mpdo_doubled_phys_trace_transfer}
    For a sector decomposition with basis elements $A_{j}$ and weights
    $\mu_{j,q}$, the physical-trace transfer of the assembled tensor is the
    block direct sum of the weighted transfers of the copies:
    \begin{align}
        \mathcal T
         & =\bigoplus_{j}\bigoplus_{q}\mu_{j,q} \mathcal T_{A_{j}}.
        \label{eq:rfp_blockwise_transfer}
    \end{align}
\end{lemma}

\begin{proof}\leanok
    The diagonal physical sum distributes over the block direct sum:
    \begin{align}
        \mathcal T
         & =\sum_i S^{ii}
        =\sum_i\bigoplus_{j,q}\mu_{j,q}A_j^{ii}
        \notag                                      \\
         & =\bigoplus_{j,q}\mu_{j,q}\sum_i A_j^{ii}
        =\bigoplus_{j,q}\mu_{j,q}\mathcal T_{A_j}.
        \notag
    \end{align}
\end{proof}

\begin{theorem}[Gauge-phase invariance of nilpotent physical trace]
    \label{thm:mpdo_doubled_phys_trace_nilpotent_gauge_phase}
    \lean{MPOTensor.%
        isNilpotent_doubledPhysTraceTransfer_of_gaugePhaseEquiv}
    \lean{MPOTensor.%
        isNilpotent_doubledPhysTraceTransfer_iff_of_gaugePhaseEquiv}
    \lean{MPOTensor.isNilpotent_doubledPhysTraceTransfer_cast_iff}
    \leanok
    \uses{def:mpdo_doubled_phys_trace_transfer, def:gauge_phase_equiv}
    If two doubled-index tensors are related by an invertible gauge and a
    nonzero phase, then the physical-trace transfer of one is nilpotent if and
    only if that of the other is nilpotent. Identifying equal bond dimensions
    does not alter this condition.
\end{theorem}

\begin{proof}\leanok
    The physical-trace transfer transforms by the same nonzero phase and
    simultaneous similarity. Powers therefore transform by the corresponding
    nonzero scalar power and similarity, which preserves whether a power
    vanishes.
\end{proof}

\begin{theorem}[Presentation independence of BNT nonnilpotency]
    \label{thm:mpdo_bnt_nonnilpotency_independent}
    \lean{MPOTensor.bnt_basis_not_isNilpotent_iff}
    \leanok
    \uses{def:cpsv_bnt_sector_presentation, def:mpdo_doubled_phys_trace_transfer}
    Let $P$ and $Q$ be BNT sector presentations of the same
    doubled-index tensor. Every representative of $P$ has non-nilpotent
    physical-trace transfer if and only if every representative of $Q$ does.
\end{theorem}

\begin{proof}\leanok
    \uses{thm:cpsv_bnt_presentation_uniqueness, thm:mpdo_doubled_phys_trace_nilpotent_gauge_phase}
    Match the representatives by
    Theorem~\ref{thm:cpsv_bnt_presentation_uniqueness}, then apply
    Theorem~\ref{thm:mpdo_doubled_phys_trace_nilpotent_gauge_phase} to each
    matched pair.
\end{proof}

\begin{definition}[Simple tensor]
    \label{def:mpdo_simple_tensor}
    \lean{MPOTensor.IsSimple}
    \leanok
    \uses{def:mpdo, def:mpdo_physical_blocking, def:cpsv_bnt_sector_presentation, def:mpdo_doubled_phys_trace_transfer}
    A tensor generating MPDOs is \emph{simple} if there is a positive
    physical blocking length $L$ such that the doubled-index tensor of the
    $L$-site blocking has a BNT sector presentation $P$ for which every
    representative $B_j$ has non-nilpotent physical-trace transfer
    $\mathcal T_{B_j}$. This records the standing canonical-block
    convention of
    \cite[lines~217--246]{Cirac2016MPDO_arXiv} together with the simplicity
    condition of \cite[lines~815--822]{Cirac2016MPDO_arXiv}.

    Every representative in $P$ has a positive number of copies, and every
    copy has nonzero weight. The blocking
    length is existential, and presentation independence follows from
    Theorem~\ref{thm:mpdo_bnt_nonnilpotency_independent}. The definition
    omits the global unit-weight normalization of
    \cite[line~246]{Cirac2016MPDO_arXiv} and does not require a normalized
    horizontal canonical-form witness.
\end{definition}

\begin{theorem}[Simplicity implies positive-length nontriviality]
    \label{thm:mpdo_simple_exists_nonzero}
    \lean{MPOTensor.IsSimple.exists_mpo_ne_zero}
    \leanok
    \uses{def:mpdo_simple_tensor, def:mpo_family}
    If $M$ is simple, then $\rho^{(N)}(M)\ne0$ for some $N>0$.
\end{theorem}

\begin{proof}\leanok
    \uses{def:cpsv_bnt_sector_presentation, lem:mpdo_three_first_site_contraction_coefficients, thm:mpdo_closed_chain_physical_blocking}
    Choose a representative and one of its copies. Every copy weight over that
    representative is nonzero, so its power-sum coefficient cannot vanish at
    every sufficiently large length. At a large length where this coefficient
    is nonzero, eventual
    linear independence of the BNT representatives makes the blocked matrix
    product vector nonzero. The closed-chain blocking identity then gives a
    nonzero unblocked MPO at a positive length.
\end{proof}

\begin{lemma}[Scaling law for MPO words]
    \label{lem:mpdo_eval_word_scalar_scaling}
    \lean{MPOTensor.evalWord_smul}
    \leanok
    \uses{def:mpo_eval_word}
    Let $M$ be an MPO tensor and let $c\in\mathbb C$. If two physical words
    $u,v$ have the same length, then
    \begin{align}
        (cM)^{u,v} & =c^{|u|}M^{u,v}.
        \notag
    \end{align}
\end{lemma}

\begin{proof}\leanok
    \uses{def:mpo_eval_word}
    Induct on the common word length and collect one factor of $c$ from each
    tensor letter.
\end{proof}

\begin{theorem}[Scaling a BNT presentation]
    \label{thm:mpdo_cpsv_bnt_scalar_scaling}
    \lean{MPSTensor.IsBNTSectorPresentation.smul_left}
    \lean{MPSTensor.SectorDecomposition.scaleWeights}
    \lean{MPSTensor.SectorDecomposition.scaleWeights_basisCount,
        MPSTensor.SectorDecomposition.scaleWeights_basisDim,
        MPSTensor.SectorDecomposition.scaleWeights_basis}
    \lean{MPSTensor.SectorDecomposition.scaleWeights_copies,
        MPSTensor.SectorDecomposition.scaleWeights_weight,
        MPSTensor.SectorDecomposition.scaleWeights_toTensor}
    \leanok
    \uses{def:cpsv_bnt_sector_presentation}
    Let $P$ be a BNT sector presentation of $A$ and let
    $c\in\mathbb C\setminus\{0\}$. Multiplying every copy weight of $P$ by
    $c$ gives a BNT sector presentation of $cA$ with the same
    representatives and multiplicities.
\end{theorem}

\begin{proof}\leanok
    \uses{lem:mpv_smul, lem:mpv_to_tensor_from_blocks_weight_mul_left}
    The new copy weights remain nonzero. Normality and eventual linear
    independence of the representatives are unchanged, and both total matrix
    product vectors acquire the common factor $c^N$ at length $N$.
\end{proof}

\begin{theorem}[Transport of a normal basis along equal MPV families]
    \label{thm:mpdo_cpsv_bnt_of_same_mpv}
    \lean{MPSTensor.IsCPSVBasisOfNormalTensors.of_sameMPV₂Pos}
    \leanok
    \uses{def:bnt_cpsv, def:same_mpv2_pos}
    If $A$ and $A'$ generate the same positive-length MPV family, then every
    basis of normal tensors for $A$ is also a basis of normal tensors for
    $A'$.
\end{theorem}

\begin{proof}\leanok
    \uses{def:bnt_cpsv, def:same_mpv2_pos}
    The basis tensors, their normality, and their eventual linear independence
    are unchanged. For every $N>0$,
    \begin{align}
        \ket{V^{(N)}(A')}
        =\ket{V^{(N)}(A)}
        =\sum_{j=1}^{g} c_j^{(N)}\ket{V^{(N)}(B_j)}.
        \notag
    \end{align}
    The first equality is Definition~\ref{def:same_mpv2_pos}, and the second is
    the spanning identity in Definition~\ref{def:bnt_cpsv}. Thus the unchanged
    coefficients $c_j^{(N)}$ give the required expansion for $A'$.
\end{proof}

\begin{theorem}[Scaling law for closed matrix-product operators]
    \label{thm:mpdo_mpo_scalar_scaling}
    \lean{MPOTensor.mpo_smul}
    \leanok
    \uses{def:mpo_family}
    Let $M$ be an MPO tensor and let $c\in\mathbb C$. For every chain length
    $N$, the closed operator generated by $cM$ is
    \begin{align}
        \rho^{(N)}(cM) & =c^N\rho^{(N)}(M).
        \notag
    \end{align}
\end{theorem}

\begin{proof}\leanok
    \uses{def:mpo_family, lem:mpdo_eval_word_scalar_scaling}
    Apply Lemma~\ref{lem:mpdo_eval_word_scalar_scaling} to the two length-$N$
    physical words, unfold the closed MPO family, and use linearity of the
    virtual trace.
\end{proof}

\begin{theorem}[Normalization is invariant under nonzero rescaling]
    \label{thm:mpdo_normalized_mpo_scaling_invariance}
    \lean{MPOTensor.normalizedMPO_smul}
    \leanok
    \uses{def:mpo_family}
    Let $M$ be an MPO tensor and let $c\in\mathbb C$ be nonzero. For every
    chain length $N$, the normalized density operators agree:
    \begin{align}
        \tr\!\left[\rho^{(N)}(cM)\right]^{-1}\rho^{(N)}(cM)
         & =\tr\!\left[\rho^{(N)}(M)\right]^{-1}\rho^{(N)}(M),
        \notag
    \end{align}
    with the inverse of zero understood to be zero.
\end{theorem}

\begin{proof}\leanok
    \uses{thm:mpdo_mpo_scalar_scaling}
    By Theorem~\ref{thm:mpdo_mpo_scalar_scaling}, the closed operator and its
    trace are both multiplied by $c^N$, and the two nonzero factors cancel in
    the normalization. When the trace vanishes, both sides are zero by the
    convention on the inverse of zero.
\end{proof}

\begin{definition}[Word encoded by a physical-block index]
    \label{def:mpdo_word_of_block}
    \lean{MPSTensor.wordOfBlock}
    \leanok
    A physical index for an $L$-site block determines a word by decoding the
    index into its $L$ original physical letters and listing them in site order.
\end{definition}

\begin{lemma}[Length of an encoded physical-block word]
    \label{lem:mpdo_length_word_of_block}
    \lean{MPSTensor.length_wordOfBlock}
    \leanok
    \uses{def:mpdo_word_of_block}
    The word encoded by an $L$-site physical-block index has length $L$.
\end{lemma}

\begin{proof}\leanok
    \uses{def:mpdo_word_of_block}
    This follows from the fixed-length block-word encoding.
\end{proof}

\begin{theorem}[Scaling law for physically blocked MPO tensors]
    \label{thm:mpdo_block_tensor_scalar_scaling}
    \lean{MPOTensor.blockTensor_smul}
    \leanok
    \uses{def:mpdo_physical_blocking}
    Let $M$ be an MPO tensor and let $c\in\mathbb C$. For every blocking length
    $L$, one has
    \begin{align}
        (cM)^{[L]} & =c^L M^{[L]}.
        \notag
    \end{align}
\end{theorem}

\begin{proof}\leanok
    \uses{def:mpdo_physical_blocking, lem:mpdo_eval_word_scalar_scaling, lem:mpdo_length_word_of_block}
    Unfold physical blocking and apply
    Lemma~\ref{lem:mpdo_eval_word_scalar_scaling} to the two encoded words;
    Lemma~\ref{lem:mpdo_length_word_of_block} gives their length $L$.
\end{proof}

\begin{theorem}[Nonnegative rescaling preserves the MPDO condition]
    \label{thm:mpdo_nonnegative_rescaling}
    \lean{MPOTensor.IsMPDO.smul_ofReal}
    \leanok
    \uses{def:mpdo}
    Let $M$ generate positive semidefinite closed operators at every positive
    length, and let $r\geq0$ be real. Then $rM$ has the same property.
\end{theorem}

\begin{proof}\leanok
    \uses{thm:mpdo_mpo_scalar_scaling}
    The length-$N$ operator is multiplied by the nonnegative scalar $r^N$.
\end{proof}

\begin{theorem}[Positive rescaling invariance of the MPDO condition]
    \label{thm:mpdo_positive_rescaling_iff}
    \lean{MPOTensor.isMPDO_smul_ofReal_iff}
    \leanok
    \uses{def:mpdo}
    Let $M$ be an MPO tensor and let $r>0$ be real. Then $rM$ generates
    positive semidefinite closed operators at every positive chain length if
    and only if $M$ does.
\end{theorem}

\begin{proof}\leanok
    \uses{thm:mpdo_nonnegative_rescaling}
    Apply Theorem~\ref{thm:mpdo_nonnegative_rescaling} first with $r$ and then
    with $r^{-1}$.
\end{proof}

\begin{theorem}[Positive rescaling preserves simplicity]
    \label{thm:mpdo_simple_positive_rescaling}
    \lean{MPOTensor.IsSimple.smul_ofReal}
    \leanok
    \uses{def:mpdo_simple_tensor}
    Let $M$ be simple and let $r>0$ be real. Then $rM$ is simple.
\end{theorem}

\begin{proof}\leanok
    \uses{thm:mpdo_nonnegative_rescaling, thm:mpdo_block_tensor_scalar_scaling, thm:mpdo_cpsv_bnt_scalar_scaling}
    Positive rescaling preserves the MPDO condition. For a blocking length
    $L$, it multiplies the blocked tensor by the nonzero scalar $r^L$.
    Multiply every copy weight by $r^L$ and apply
    Theorem~\ref{thm:mpdo_cpsv_bnt_scalar_scaling}; the representatives and
    their non-nilpotent physical-trace transfers are unchanged.
\end{proof}

\begin{theorem}[Positive rescaling invariance of simplicity]
    \label{thm:mpdo_simple_positive_rescaling_iff}
    \lean{MPOTensor.isSimple_smul_ofReal_iff}
    \leanok
    \uses{def:mpdo_simple_tensor}
    Let $M$ be an MPO tensor and let $r>0$ be real. Then $rM$ is simple
    if and only if $M$ is simple.
\end{theorem}

\begin{proof}\leanok
    \uses{thm:mpdo_simple_positive_rescaling}
    Apply the preceding implication first with $r$ and then with $r^{-1}$.
\end{proof}
\begin{definition}[Normalized simplicity in blocked canonical form]
    \label{def:mpdo_simple_cf_tensor}
    \lean{MPOTensor.IsSimpleCanonicalForm}
    \leanok
    \uses{def:mpdo, def:horizontal_cf_mpdo, def:mpdo_doubled_phys_trace_transfer}
    A tensor generating MPDOs satisfies this predicate if it has a normalized
    horizontal basis-of-normal-tensors decomposition whose basis elements all
    have non-nilpotent physical-trace transfer. This records the
    horizontal-canonical-form data of the already-blocked tensor $K$ fixed
    in~\cite[Appendix~C.2, line~1628]{Cirac2016MPDO_arXiv}.

    The definition concerns this fixed representative. Its witness includes
    the global unit-weight convention
    of~\cite[line~246]{Cirac2016MPDO_arXiv}, whereas the non-nilpotency
    condition
    itself is unchanged by nonzero scalar rescaling. It supplies the
    Appendix~C.2 hypothesis on the already-blocked tensor and is not the
    simplicity predicate of Definition~\ref{def:mpdo_simple_tensor},
    which quantifies over blocking lengths and presentations. No invariance
    under rescaling or replacement of the canonical presentation is asserted.
\end{definition}
% Note for maintainers: this predicate does not include the biCF
% (bidirectional canonical form) hypothesis imposed on K at line 1628.  That
% condition is not used by the copy-independence argument and is introduced
% separately when the inverse tensor is needed.

\begin{lemma}[Normalized simple canonical forms are horizontal canonical]
    \label{lem:mpdo_simple_tensor_horizontal_cf}
    \lean{MPOTensor.IsSimpleCanonicalForm.isHorizontalCF}
    \leanok
    \uses{def:mpdo_simple_cf_tensor, def:horizontal_cf_mpdo}
    A tensor satisfying normalized simplicity in the chosen blocked
    canonical-form setting is in horizontal canonical form.
\end{lemma}

\begin{proof}\leanok
    The canonical-form data in the simplicity hypothesis---the sector
    decomposition, its basis-of-normal-tensors property, and the
    block-diagonal gauge---give a horizontal canonical form. Positivity and
    non-nilpotency are not needed for this conclusion.
\end{proof}
